% Сборка: pdflatex или xelatex -interaction=nonstopmode mathematical_model_and_mesh.tex
% Внимание: текущий точечный генератор НЕ является точным отображением
% области с конечными скважинами. Корректная постановка дана в разделе 6.
\documentclass[12pt,a4paper]{article}
\usepackage{iftex}
\ifPDFTeX
  \usepackage[T2A]{fontenc}
  \usepackage[utf8]{inputenc}
\else
  \usepackage{fontspec}
  \setmainfont{Times New Roman}
\fi
\usepackage[english,russian]{babel}
\usepackage{amsmath,amssymb,mathtools}
\usepackage{geometry}
\usepackage{booktabs}
\usepackage{enumitem}
\usepackage{hyperref}
\geometry{left=25mm,right=20mm,top=22mm,bottom=23mm}
\hypersetup{hidelinks,pdftitle={Математическая постановка и построение сетки PorousFlow 2D}}
\setlist{nosep}
\allowdisplaybreaks

\title{Математическая постановка и реализация построения сетки\\
       для двумерной ячейки пятиточечной системы}
\author{Техническое описание реализации PorousFlow 2D}
\date{17 сентября 2026 г.}

\begin{document}
\maketitle
\tableofcontents
\newpage

\section{Назначение и границы модели}

Рассматривается вытеснение нефти водой в жёстком пористом пласте постоянной
толщины $h$. Обе жидкости считаются несжимаемыми и несмешивающимися.
Капиллярное давление и гравитация в текущем расчёте отключены. Неизвестные
поля --- общее давление $p(x,y,t)$ и водонасыщенность $S(x,y,t)$. Ниже сначала
задана физическая задача, затем её конечно-объёмная аппроксимация и два
реализованных способа построения структурированной сетки. Основное внимание
уделено сопоставлению текущего приближённого генератора с требуемым
конформным отображением в координатах потенциала $\varphi$ и функции тока
$\psi$. Текущий генератор \emph{не} решает в точности задачу с конечными
радиусами скважин; это существенное ограничение, а не погрешность
визуализации.

Следует различать \emph{вспомогательное аналитическое поле}, используемое
для построения сетки, и \emph{фактическое поле течения}, которое решается на
этой сетке. При неоднородной проницаемости или при изменении подвижности из-за
вытеснения они вообще не совпадают.

\section{Расчётная область и условия на её границе}

Пусть $L_x,L_y>0$, а центры нагнетательной и добывающей скважин находятся в
противоположных вершинах прямоугольника:
\begin{equation}
 \mathbf x_I=(0,0),\qquad \mathbf x_P=(L_x,L_y).
\end{equation}
Из прямоугольника вырезаны четверти дисков радиусов $r_I$ и $r_P$:
\begin{equation}
 \Omega=\bigl([0,L_x]\times[0,L_y]\bigr)
 \setminus\bigl(B_{r_I}(\mathbf x_I)\cup B_{r_P}(\mathbf x_P)\bigr).
 \label{eq:domain}
\end{equation}
Дуги скважин обозначаются $\Gamma_I$ и $\Gamma_P$. Внешняя непроницаемая
граница разбита на два пути: $\Gamma_A$ идёт от нагнетателя по нижней и правой
сторонам, $\Gamma_B$ --- по левой и верхней сторонам. Свободные углы
$(L_x,0)$ и $(0,L_y)$ являются точками стагнации вспомогательного
аналитического течения. Четверти скважин при зеркальном продолжении ячейки
образуют полные скважины шахматной пятиточечной системы.

В реализованном сценарии на скважинах заданы давления, а на прямых сторонах
--- отсутствие потока:
\begin{align}
 p|_{\Gamma_I}&=p_I, & p|_{\Gamma_P}&=p_P, & p_I&>p_P,\label{eq:bc-pressure}\\
 \mathbf u_t\cdot\mathbf n&=0 &&\text{на }\Gamma_A\cup\Gamma_B.
 \label{eq:bc-noflow}
\end{align}
На входящем потоке нагнетателя задаётся доля воды $f_w^I\in[0,1]$;
по умолчанию $f_w^I=1$. На добывающей скважине применяется адвективный
выход: состав продукции определяется насыщенностью приграничной ячейки.
Начальное условие имеет вид $S(x,y,0)=S_0(x,y)$; пример проекта использует
постоянное $S_0$.

\section{Уравнения двухфазной фильтрации и переноса}

Для фаз $\alpha\in\{w,o\}$ закон Дарси и подвижности имеют вид
\begin{equation}
 \mathbf u_\alpha=-\mathbf K\lambda_\alpha(S)\nabla p,
 \qquad
 \lambda_\alpha(S)=\frac{k_{r\alpha}(S)}{\mu_\alpha},
 \qquad S_o=1-S.
 \label{eq:darcy}
\end{equation}
Здесь $\mathbf K(x,y)$ --- абсолютная проницаемость (скаляр или симметричный
положительно определённый тензор), $\mu_\alpha$ --- динамическая вязкость,
$k_{r\alpha}$ --- относительная фазовая проницаемость. Полный поток и
дробная функция воды определяются как
\begin{equation}
 \lambda_t=\lambda_w+\lambda_o,\qquad
 \mathbf u_t=-\mathbf K\lambda_t(S)\nabla p,\qquad
 f_w(S)=\frac{\lambda_w(S)}{\lambda_t(S)}.
 \label{eq:total-flux}
\end{equation}
При отсутствии распределённых внутренних источников система уравнений в
$\Omega$ записывается в виде
\begin{align}
 \nabla\cdot\mathbf u_t&=0,
 \label{eq:pressure-pde}\\
 \phi_p(x,y)\,\frac{\partial S}{\partial t}
 +\nabla\cdot\bigl(f_w(S)\mathbf u_t\bigr)&=0.
 \label{eq:transport-pde}
\end{align}
$\phi_p$ обозначает пористость; индекс нужен, чтобы не путать её с
аналитическим потенциалом $\varphi$. Вода и нефть поступают в область или
выходят из неё через граничные дуги скважин, поэтому объёмные источники в
\eqref{eq:pressure-pde}--\eqref{eq:transport-pde} не требуются.

Реализованы степенные кривые относительных проницаемостей. При связанной
водонасыщенности $S_{wc}$ и остаточной нефтенасыщенности $S_{or}$:
\begin{align}
 S_e&=\operatorname{clip}\!\left(
    \frac{S-S_{wc}}{1-S_{wc}-S_{or}},,0,,1\right),\\
 k_{rw}(S)&=k_{rw}^{\rm end}S_e^{n_w},\qquad
 k_{ro}(S)=k_{ro}^{\rm end}(1-S_e)^{n_o}.
 \label{eq:relperm}
\end{align}
Физический диапазон: $S_{wc}\le S\le 1-S_{or}$. Параметры
$k_{r\alpha}^{\rm end}$ и $n_\alpha$ задаются во входном проекте.

\section{Конечно-объёмная дискретизация}

Узлы сетки хранятся как $\mathbf X_{ij}\in\mathbb R^2$, где
$i=0,\ldots,N_x$ задаёт продольное направление от нагнетателя к добывающей
скважине, $j=0,\ldots,N_y$ --- поперечное. Ячейка
$C_{ij}$ ограничена четырьмя узлами
$\mathbf X_{ij},\mathbf X_{i+1,j},\mathbf X_{i+1,j+1},\mathbf X_{i,j+1}$.
Давление и насыщенность размещаются в центре ячейки. Для грани $f$ с длиной
$\ell_f$ её физическая площадь равна $A_f=h\ell_f$, а поровый объём ячейки
равен $V^p_{ij}=h\,|C_{ij}|\phi_{p,ij}$.

Проводимость внутренней грани между ячейками $L$ и $R$ вычисляется через
нормальные проницаемости $K_{n,L}=\mathbf n_f^\mathsf T\mathbf K_L\mathbf n_f$
и $K_{n,R}$. Если $d_L,d_R$ --- положительные проекции расстояний от центров
ячеек до грани на её нормаль, то
\begin{equation}
 G_L=\frac{A_f\lambda_{t,L}K_{n,L}}{d_L},\qquad
 G_R=\frac{A_f\lambda_{t,R}K_{n,R}}{d_R},\qquad
 T_f=\frac{G_LG_R}{G_L+G_R}.
 \label{eq:transmissibility}
\end{equation}
Полный расход, положительный из $L$ в $R$, равен
\begin{equation}
 Q_{t,f}=T_f(p_L-p_R).
 \label{eq:two-point}
\end{equation}
На грани с заданным давлением используется односторонняя проводимость.
Грани симметрии имеют $Q_{t,f}=0$. Сумма ориентированных гранных расходов
каждой ячейки равна нулю. Полученная симметричная положительно определённая
система для давления решается PCG с диагональным предобуславливателем;
после решения проверяется локальный баланс гранных расходов.
Формула~\eqref{eq:two-point} является двухточечной аппроксимацией;
поправка на неортогональность в текущей версии отсутствует. Поэтому для
сильно косых ячеек и произвольной полной тензорной $\mathbf K$ одной
положительности площади недостаточно для гарантии точности.

Водный расход через грань задаётся консервативно по направлению полного
расхода:
\begin{equation}
 Q_{w,f}=Q_{t,f}\,f_w(S_{\mathrm{up}(f)}).
 \label{eq:upwind}
\end{equation}
На входе нагнетателя вместо $f_w(S_{\mathrm{up}})$ подставляется $f_w^I$.
Неявное обновление насыщенности имеет вид
\begin{equation}
 \frac{V^p_{ij}}{\Delta t}(S_{ij}^{n+1}-S_{ij}^{n})
 +\sum_{f\subset\partial C_{ij}}Q_{w,f}(S^{n+1})=0,
 \label{eq:implicit}
\end{equation}
где ориентированные расходы считаются положительными наружу. Есть также
явная опция, в которой в гранных функциях используется $S^n$ и шаг
ограничивается по CFL. В обеих опциях давление и полный расход сначала
пересчитываются из текущей насыщенности, после чего выполняется перенос.
Показатели продукции определяются суммированием выходящих гранных расходов
на $\Gamma_P$:
\begin{equation}
 Q_{\rm liq}^{P}=\sum_{f\subset\Gamma_P}Q_{t,f},\quad
 Q_w^{P}=\sum_{f\subset\Gamma_P}Q_{w,f},\quad
 \mathrm{WCT}^{P}=Q_w^P/Q_{\rm liq}^P.
\end{equation}

\section{Топология сетки и два генератора}

Для обоих вариантов область трактуется как криволинейный четырёхугольник:
дуги скважин образуют стороны $i=0$ и $i=N_x$, а два непроницаемых
L-образных пути --- стороны $j=0$ и $j=N_y$. Это обеспечивает покрытие
всей области, включая две свободные угловые зоны. Геометрия узлов, нормали,
длины граней, площади и объёмы вычисляются из координат $\mathbf X_{ij}$.

\subsection{Гладкая геометрически согласованная сетка}

В варианте \texttt{body\_fitted\_structured} четыре граничные кривые
задаются точно: $\mathbf C_0(\eta)$ и $\mathbf C_1(\eta)$ --- дуги
скважин, $\mathbf D_0(\xi)$ и $\mathbf D_1(\xi)$ --- два внешних пути.
Для внутренних узлов применяется трансфинитная интерполяция
\begin{align}
 \mathbf X(\xi,\eta)={}&(1-\xi)\mathbf C_0(\eta)
 +\xi\mathbf C_1(\eta)
 +(1-\eta)\mathbf D_0(\xi)+\eta\mathbf D_1(\xi)
 \notag\\[-1mm]
 &-\bigl[(1-\xi)(1-\eta)\mathbf P_{00}
 +(1-\xi)\eta\mathbf P_{01}
 +\xi(1-\eta)\mathbf P_{10}+\xi\eta\mathbf P_{11}\bigr].
 \label{eq:tfi}
\end{align}
Продольные уровни $\xi_i$ строятся симметричной геометрической прогрессией:
мелкие ячейки у обеих скважин, более крупные в середине. Расчётная сетка
следует геометрии, но её внутренние линии \emph{не обязаны} совпадать с
фактическими линиями тока и изобарами.

\subsection{Историческая версия точечного генератора
\texorpdfstring{$(\varphi,\psi)$}{(phi, psi)}}

Далее в этом подразделе зафиксирован прежний алгоритм с круговыми границами и
дополнительным перераспределением линий. Он больше не используется программой.
Текущая постановка с изобарическими границами, равномерными уровнями функции
тока и естественным логарифмическим сгущением приведена в отдельном документе
\texttt{point\_source\_isobar\_mesh.tex}.

Вариант \texttt{analytic\_phi\_psi} \emph{в его текущей реализации} использует вспомогательное однородное
изотропное течение в зеркально продолженной пятиточечной системе. Источники
и стоки расположены в точках
\begin{equation}
 \mathbf x_{I,mn}=(2mL_x,2nL_y),\qquad
 \mathbf x_{P,mn}=((2m+1)L_x,(2n+1)L_y),\qquad m,n\in\mathbb Z.
\end{equation}
На практике бесконечная
решётка не обрезается прямоугольным окном образов, поскольку такая обрезка
создаёт небольшой, но важный около сепаратрис паразитный поток через стороны.
Вместо этого используется представление через тета-функцию.

Для устойчивой сходимости рядов координаты при необходимости меняются
местами так, чтобы $a=\min(L_x,L_y)$ был действительным полупериодом,
$b=\max(L_x,L_y)$ --- мнимым. Пусть $z=x'+\mathrm i y'$ обозначает
полученную комплексную координату и $z_P=a+\mathrm i b$. Введём
\begin{align}
 q&=\exp(-\pi b/a),\\
 \Theta_N(\zeta;q)&=\sum_{n=0}^{N-1}(-1)^n
 q^{(n+1/2)^2}\sin((2n+1)\zeta),\qquad N=7\ \text{по умолчанию}.
 \label{eq:theta}
\end{align}
Множитель $2$ стандартной функции $\theta_1$ опущен, поскольку он
сокращается в отношении. Код при вычислении использует не менее четырёх
членов ряда. Вспомогательный комплексный потенциал можно записать как
\begin{equation}
 F(z)=\log\Theta_N\!\left(\frac{\pi z}{2a};q\right)
 -\log\Theta_N\!\left(\frac{\pi(z-z_P)}{2a};q\right)
 +\frac{\mathrm i\pi z}{2a},\qquad
 \varphi=\operatorname{Re}F.
 \label{eq:complex-potential}
\end{equation}
Это потенциал \emph{точечных} источников и стоков, находящихся в центрах
вырезанных дисков. Он обеспечивает симметрию внешних сторон, но в общем
случае $\varphi|_{\Gamma_I}$ и $\varphi|_{\Gamma_P}$ не являются
константами. Поэтому~\eqref{eq:complex-potential} нельзя считать точным
решением задачи~\eqref{eq:domain}--\eqref{eq:bc-noflow} при конечных
$r_I,r_P$. Тот факт, что при очень малом $r/L$ разброс потенциала по дуге
мало заметен на рисунке, не устраняет нарушение граничного условия.

Программная реализация не вычисляет комплексный логарифм с выбором ветви:
она использует разность логарифмов модулей $\Theta_N$ для $\varphi$ и
логарифмическую производную для вектора $\mathbf v=\nabla\varphi$.
Линейное слагаемое в~\eqref{eq:complex-potential} компенсирует
квазипериодичность тета-функции по мнимому направлению. В несменённой
координатной системе
\begin{equation}
 v_x-\mathrm i v_y=\frac{\pi}{2a}
 \left[
 \frac{\Theta_N'(\pi z/2a;q)}{\Theta_N(\pi z/2a;q)}
 -\frac{\Theta_N'(\pi(z-z_P)/2a;q)}
 {\Theta_N(\pi(z-z_P)/2a;q)}+\mathrm i
 \right].
 \label{eq:complex-velocity}
\end{equation}
При перестановке координат компоненты $\mathbf v$ переставляются обратно.
У $\mathbf v$ нет заданного физического масштаба расхода: её назначение
--- построить геометрические линии. Для однородного изотропного пласта
изобары $p$ и уровни $\varphi$ совпадают с точностью до монотонного
преобразования только при соответствующих допущениях о подвижности.

\paragraph{Опорная изобара и стартовые точки.}
На квадратной ячейке потенциалы свободных углов одинаковы; опорный уровень
\begin{equation}
 \varphi_*=\frac{\varphi(L_x,0)+\varphi(0,L_y)}{2}
 \label{eq:reference-level}
\end{equation}
проходит через оба угла. На дуге нагнетателя сначала берётся густой набор
пилотных параметров $\eta\in[0,1]$ с квадратичным сгущением у концов.
Из каждой стартовой точки интегрируется линия тока до $\varphi_*$.
Пересечения образуют дискретную опорную кривую. Её нормированная накопленная
длина $s(\eta)$ вычисляется суммированием длин соседних отрезков.

Для $u_j=j/N_y$ задаются целевые положения
\begin{equation}
 G(u)=\begin{cases}
 2u^2,&0\le u\le\tfrac12,\\
 1-2(1-u)^2,&\tfrac12<u\le1,
 \end{cases}\qquad
 s_j=(1-\beta)u_j+\beta G(u_j),\quad 0\le\beta\le1.
 \label{eq:streamline-grading}
\end{equation}
Параметр $\beta$ называется в интерфейсе «Линии у краёв»
(\texttt{streamline\_edge\_fraction}). Стартовые параметры находятся
численным обращением $\eta_j=s^{-1}(s_j)$. Узлы на двух крайних линиях
$j=0,N_y$ берутся непосредственно с непроницаемых сторон области.
Именно распределение по длине \emph{всей} опорной кривой, а не равное
разбиение расхода на маленькой дуге скважины, позволяет провести линии через
угловые зоны при небольшом $N_y$.

\paragraph{Уровни изобар и сгущение у скважин.}
Пусть $L_B=L_x+L_y-r_I-r_P$ --- длина каждого L-образного внешнего пути.
Расчётный минимальный продольный масштаб согласуется с поперечным размером
ячейки на дугах скважин:
\begin{equation}
 \delta=\operatorname{clip}\!\left(
 \frac{\pi(r_I+r_P)}{4N_yL_B},10^{-7},5\cdot10^{-3}\right).
 \label{eq:grading-core}
\end{equation}
При заданном коэффициенте сгущения $B>0$ функция распределения вдоль
внешнего пути имеет вид
\begin{equation}
 M(s)=1+\frac{0.02B}{s+\delta}
       +\frac{0.02B}{1-s+\delta},\qquad 0\le s\le1.
 \label{eq:monitor}
\end{equation}
Уровни $s_i$ выбираются как квантили накопленного интеграла $M$:
\begin{equation}
 \frac{\int_0^{s_i}M(\sigma)\,d\sigma}
      {\int_0^1M(\sigma)\,d\sigma}=\frac{i}{N_x},
 \qquad i=0,\ldots,N_x.
 \label{eq:monitor-inverse}
\end{equation}
В коде интеграл и его обратная функция вычисляются на густой табличной
сетке. Обе внешние границы параметризованы нормированной длиной.
Общий целевой уровень определяется средним двух значений потенциала:
\begin{equation}
 \widetilde\varphi_i=\tfrac12\left[
 \varphi(\mathbf D_0(s_i))+\varphi(\mathbf D_1(s_i))\right].
 \label{eq:phi-target}
\end{equation}
После ограничения общим диапазоном потенциала выполняются монотонизация и
малая строго возрастающая регуляризация порядка $10^{-10}$. Уровень $B=1$
даёт обычное сгущение, $B=2$ --- сильное, $B=5$ --- очень сильное.
Постоянная единица в~\eqref{eq:monitor} не даёт убрать все изобары из центра.

\paragraph{Трассировка и сборка узлов.}
Поскольку $\mathbf v=\nabla\varphi$, вдоль линии $\psi=\mathrm{const}$
можно использовать $\varphi$ как независимую переменную:
\begin{equation}
 \frac{d\mathbf x}{d\varphi}
 =\frac{\mathbf v(\mathbf x)}{|\mathbf v(\mathbf x)|^2}.
 \label{eq:trace-ode}
\end{equation}
В реализации применяется RK4 с адаптивным шагом: один полный шаг
сравнивается с двумя половинными, ошибка управляет шагом, а узлы выдаются
на общих уровнях $\varphi_i$. Для внешних линий пересечения с уровнями
$\varphi_i$ находятся интерполяцией по точно заданной границе. Первый ряд
узлов фиксируется на дуге нагнетателя; последний проектируется на дугу
добывающей скважины. Таким образом, узлы имеют логическую структуру
$\mathbf X_{ij}\approx\mathbf x(\varphi_i,\psi_j)$, но концам у скважин
приоритетно обеспечивается геометрическая точность.

\paragraph{Ограничения текущего генератора.}
При $L_x\ne L_y$ значения $\varphi$ в свободных углах вообще различны.
Текущий алгоритм всё равно использует среднее~\eqref{eq:reference-level}
и принудительно фиксирует крайние точки опорной кривой в углах. Поэтому
для прямоугольной ячейки опорная кривая в самых концах не является точной
изобарой. Кроме того, в точках стагнации $|\mathbf v|=0$, и
\eqref{eq:trace-ode} сингулярно. Сильная градация размеров у скважин и
вытянутые ячейки возможны даже при положительных площадях и хороших углах.
Даже для квадратной однородной изотропной ячейки такая сетка остаётся
приближением к задаче с конечными скважинами. Численное обращение
\eqref{eq:complex-potential} вместо трассировки~\eqref{eq:trace-ode}
\emph{не исправит} этого нарушения: исходное комплексное отображение
сначала необходимо заменить.

\section{Требуемое отображение ТФКП для скважин конечного радиуса}

Правильный объект для этой геометрии --- конформное отображение
\emph{самой области}~\eqref{eq:domain}, а не потенциал особенностей,
поставленных в центры скважин. Сначала определяется единственная
гармоническая функция $u$ из смешанной краевой задачи
\begin{equation}
 \begin{cases}
  \Delta u=0,& z\in\Omega,\\
  u=0,&z\in\Gamma_I,\\
  u=1,&z\in\Gamma_P,\\
  \partial_nu=0,&z\in\Gamma_A\cup\Gamma_B.
 \end{cases}
 \label{eq:exact-map-bvp}
\end{equation}
Здесь радиусы $r_I,r_P$ входят непосредственно в геометрию границ, а
равенства $u=0$ и $u=1$ должны выполняться на \emph{всех точках} двух
круговых дуг, не только на их концах. Для однородного изотропного пласта
с постоянной подвижностью физическое давление равно
$p=p_I+(p_P-p_I)u$.

Гармонически сопряжённая функция $v$ определяется уравнениями
Коши--Римана
\begin{equation}
 v_x=-u_y,\qquad v_y=u_x.
 \label{eq:cr}
\end{equation}
Положив $v=0$ на $\Gamma_A$, получаем $v=M$ на $\Gamma_B$, где
\begin{equation}
 M=-\int_{\Gamma_I}\partial_nu\,ds
  =\int_{\Gamma_P}\partial_nu\,ds>0.
 \label{eq:modulus}
\end{equation}
Тогда $W(z)=u(z)+\mathrm i v(z)$ отображает криволинейный
четырёхугольник $\Omega$ на прямоугольник
\begin{equation}
 R_M=\{w\in\mathbb C:0<\operatorname{Re}w<1,\;
                         0<\operatorname{Im}w<M\}.
 \label{eq:canonical-rectangle}
\end{equation}
В частности, обе скважинные дуги целиком переходят в две вертикальные
стороны, а оба внешних непроницаемых пути --- в горизонтальные стороны.
Такое представление через смешанную задачу и сопряжённую функцию является
стандартным методом определения конформного модуля четырёхугольника
\cite{conjugate-map}. Для границ с круговыми дугами возможно также
вычисление отображения методами обобщённого преобразования
Шварца--Кристоффеля; его дополнительные параметры обычно определяются
численно \cite{circular-arc-map}.

\subsection{Узлы действительно конформной сетки}

После нахождения $W$ выбираются строго возрастающие последовательности
$U_i\in[0,1]$ и $V_j\in[0,M]$ с желаемым сгущением. Тогда
\begin{equation}
 \boxed{\quad z_{ij}=W^{-1}(U_i+\mathrm i V_j),\qquad
 i=0,\ldots,N_x,\quad j=0,\ldots,N_y.\quad}
 \label{eq:conformal-nodes}
\end{equation}
Линии $U_i=\mathrm{const}$ являются изобарами, линии
$V_j=\mathrm{const}$ --- линиями тока для \emph{этой} краевой задачи.
Точно заданные дуги скважин сохраняются как границы, без радиальной
проекции узлов после вычисления сетки. Численное обращение уже
\emph{правильного} отображения допустимо выполнять методом Ньютона с
продолжением по соседним узлам:
\begin{equation}
 z^{(k+1)}=z^{(k)}-
 \frac{W(z^{(k)})-(U_i+\mathrm iV_j)}{W'(z^{(k)})}.
 \label{eq:inverse-map}
\end{equation}
Это численное решение уравнения для узла, а не численная трассировка
линии тока. Однако построение самого $W$ для данного контура с дугами
обычно требует численного определения параметров конформного отображения;
обещать замкнутую элементарную формулу для произвольных $L_x,L_y,r_I,r_P$
было бы необоснованно.

\subsection{Критерии корректности и статус реализации}

Проверять нужно не только положительные площади ячеек, но и невязки
краевых условий на \emph{всех} дугах и прямых сторонах:
\begin{equation}
 \max_{\Gamma_I}|u|,\quad
 \max_{\Gamma_P}|u-1|,\quad
 \max_{\Gamma_A\cup\Gamma_B}|\partial_nu|,
 \qquad
 \max_{\Omega}|\nabla v-J\nabla u|,
 \quad J(a,b)=(-b,a).
 \label{eq:map-checks}
\end{equation}
Дополнительно проверяются взаимная ортогональность внутренних координатных
линий, отсутствие пересечений и сходимость координат узлов при повышении
точности вычисления $W$.

\textbf{Статус:} отображение~\eqref{eq:exact-map-bvp}--\eqref{eq:conformal-nodes}
\emph{ещё не реализовано} в текущей версии программы. Имеющийся режим
\texttt{analytic\_phi\_psi} основан на точечных особенностях и не должен
выдаваться за точное построение сетки для скважин конечного радиуса.

\section{Проверки и связь с программой}

Перед расчётом программа проверяет знак площади каждого четырёхугольника
и минимальный внутренний угол. Порог минимального угла составляет $40^\circ$
для гладкого генератора и $2^\circ$ для аналитического. Максимальное
отношение длин сторон выводится как диагностика, но для аналитического
варианта жёстко не ограничивается. Проверка по углу и знаку площади
\emph{не доказывает} достаточности сетки для достоверного решения переноса;
необходим также анализ сходимости при измельчении сетки и баланса масс.

\begin{center}
\begin{tabular}{@{}ll@{}}
\toprule
Компонент & Файл реализации \\
\midrule
Геометрия и трансфинитная интерполяция & \texttt{src/porous\_fvm/geometry.py} \\
Аналитический потенциал и трассировка & \texttt{src/porous\_fvm/analytic\_mesh.py} \\
Гладкое распределение узлов & \texttt{src/porous\_fvm/flow\_mesh.py} \\
Геометрические метрики и контроль качества & \texttt{src/porous\_fvm/mesh.py} \\
Давление и гранные расходы & \texttt{src/porous\_fvm/flow.py} \\
Перенос насыщенности & \texttt{src/porous\_fvm/transport.py} \\
\bottomrule
\end{tabular}
\end{center}

Соответствие математических и входных параметров:
\begin{itemize}
 \item $N_x$: \texttt{mesh.cells.along\_streamline};
 \item $N_y$: \texttt{mesh.cells.across\_streamlines};
 \item способ построения: \texttt{mesh.generator};
 \item $B$: \texttt{mesh.pressure\_bias};
 \item $\beta$: \texttt{mesh.streamline\_edge\_fraction}.
\end{itemize}
Изменение этих параметров требует повторного построения сетки до запуска
расчёта.

\begin{thebibliography}{9}
\bibitem{conjugate-map}
H.~Hakula, A.~Rasila.
\emph{Laplace--Beltrami Equations and Numerical Conformal Mappings on Surfaces}.
SIAM Journal on Scientific Computing, 2025.
\url{https://doi.org/10.1137/24M1656840}.

\bibitem{circular-arc-map}
P.~Bjørstad, E.~Grosse.
\emph{Conformal Mapping of Circular Arc Polygons}.
SIAM Journal on Scientific and Statistical Computing.
\url{https://doi.org/10.1137/0908003}.
\end{thebibliography}

\end{document}
